\documentclass[11pt]{article}

\usepackage[a4paper,margin=30mm]{geometry}
\usepackage{amsmath,amssymb,amsthm,mathtools}
\usepackage[hidelinks]{hyperref}
\usepackage{microtype}

\newtheorem{theorem}{Theorem}[section]
\newtheorem{corollary}[theorem]{Corollary}
\newtheorem{lemma}[theorem]{Lemma}
\newtheorem{remark}[theorem]{Remark}

\newcommand{\R}{\mathbb R}
\newcommand{\N}{\mathbb N}
\newcommand{\Lam}{\Lambda}
\newcommand{\A}{\mathcal A}
\newcommand{\w}{\mathsf w}
\newcommand{\Sur}{\mathsf S}

\title{Sharp Tilted-Tent Laws and Unbounded\\Suzuki--Mangoldt Prime-Block Surplus}
\author{HGT Research Project}
\date{22 July 2026}

\begin{document}
\maketitle

\begin{abstract}
We establish two exact laws for exponentially tilted tent profiles and use them to solve the
source-scaling problem at aligned Suzuki--Mangoldt apices.  First, the lower envelope of two
geometric-bracket endpoint profiles has an explicit sharp minimax constant.  Second, the normalized
right flank
\[
 R_{\beta,L}(d)=e^{\beta d}(1-d/L)
\]
is strictly decreasing whenever $\beta L<1$; more strongly, $R_{\beta,L}(r)$ is the exact least
uniform ceiling on every tail $[r,L]$.

At width $L=\log5$, let the moving Suzuki apex be the Mangoldt node $2\log n$, and let
$P_n=\{p\text{ prime}:2n<p\le5n\}$.  The continuous tail law gives the sharp contraction
\[
 q_5=\sqrt2\left(1-\frac{\log2}{\log5}\right)\in(0,1)
\]
for every $p\in P_n$.  A symbolic Chebyshev argument proves, for $n\ge e^{20}$,
\[
 \sum_{p\in P_n}\frac{2\log p}{p}\ge\frac1{10}.
\]
Consequently the actual-prime source surplus obeys the macroscopic law
\[
 \Sur_5(n)\ge
 \frac{1-q_5}{10}\sqrt n\log5>0.
\]
In particular, $\Sur_5$ is unbounded.  This replaces the vanishing $O(n^{-1/2})$ single-flank
certificate by an unconditional $\Omega(\sqrt n)$ block certificate, the scale of a fixed-width
moving-apex seam.  We retain the sharp dyadic law as the local extremal model.  Every mathematical
claim is checked in Lean~4 without RH, a prime-gap hypothesis, a capacity premise, or numerical
certification.  The remaining obstruction is target-side capacity domination, not source scale.
\end{abstract}

\section{Introduction}\leavevmode\par

Suzuki's Hilbert-space program packages the Riemann hypothesis into positivity and extension
questions for concrete spaces and operators associated with the completed zeta function
\cite{Suzuki2026}.  In the rectangular lane, exact Fourier and autocorrelation identities convert
the prime contribution into a positive lag channel sampled at the logarithmic Mangoldt nodes.
After monotone branches are transported, the remaining obstruction occurs at a moving maximum:
an individual atom at the maximum cannot be paid by a positive-measure pointwise superlevel set.

Two neighboring literatures delimit the contribution.  Groskin's finite Guinand--Weil dictionary
identifies exact finite quadratic forms and controls an omitted archimedean tail, while making no
prime-counting claim \cite{Groskin2026}.  Transport-majorization supplies general convex-order
machinery from controlled transport maps \cite{MelbourneRoberto2021}.  Neither scope supplies the
source-specific object required here: an actual Mangoldt block aligned to a Suzuki moving apex
with a sharp contraction and macroscopic surplus.  The result below constructs that object; it
does not claim the still-missing target-capacity comparison.

The natural next object is therefore a block, not an individual atom.  A block must contain the
apex atom together with lower-valued source atoms, so its weighted average is strictly below the
apex.  Two issues are logically distinct:
\begin{enumerate}
\item produce a source-canonical block with a proved positive sample deficit;
\item construct disjoint continuum target blocks whose kernel excess dominates that deficit.
\end{enumerate}
Earlier local certificates settle strict positivity but decay like $n^{-1/2}$.  A fixed-width
future/past overlap at a moving apex has height proportional to $\sqrt n$, so positivity alone does
not answer the scale question.  This paper settles that question: an explicit source-canonical
prime block produces surplus proportional to $\sqrt n$, with no prime-gap input.  It makes no
assertion about target domination and no assertion of RH.

There are four complementary laws.  The cross-prime three-atom theorem proves local dilution.
The geometric-bracket theorem gives the exact two-endpoint minimax constant.  The subcritical
tail theorem gives an exact one-sided envelope and, at width $\log5$, converts a macroscopic prime
block into a surplus of order $\sqrt n$.  The dyadic boundary identity remains the sharp local
model: an adjacent zero-sample boundary atom is useful for dilution even though its mass alone is
too small to control the macroscopic seam.

All mathematical statements below are mirrored by named declarations in four Lean modules.  The
formal theorem map is recorded in Section~\ref{sec:formal}.

\section{Aligned samples and exact support}\leavevmode\par

For $n\ge2$, set
\[
 a_n=\frac12\log(2n),\qquad \tau_m=2\log m.
\]
Let $\A_n(m)$ denote the width-$\log 2$ positive rectangular Suzuki channel at aperture $a_n$,
sampled at $\tau_m$.  The previously established exact sample formula is
\begin{equation}\label{eq:sample}
 \A_n(m)=\sqrt m\,
 \max\!\left\{0,\log 2-\left|\log\frac mn\right|\right\}.
\end{equation}
The Mangoldt mass attached to the node is
\begin{equation}\label{eq:mass}
 \w_m=\frac{2\Lam(m)}m.
\end{equation}

\begin{theorem}[Exact open-annulus support]\label{thm:support}
For all integers $n,m\ge2$,
\[
 \A_n(m)>0
 \quad\Longleftrightarrow\quad
 n<2m\ \text{ and }\ m<2n.
\]
\end{theorem}

\begin{proof}
The factor $\sqrt m$ in \eqref{eq:sample} is positive.  The remaining factor is positive exactly
when
\[
 \left|\log\frac mn\right|<\log2.
\]\leavevmode\par
Strict monotonicity of the logarithm on $(0,\infty)$ converts this inequality to
$1/2<m/n<2$, which is the claimed pair of integer inequalities.
\end{proof}

At the apex,
\begin{equation}\label{eq:self}
 \A_n(n)=\sqrt n\log2>0.
\end{equation}
For $n\ge4$, the aligned node is the unique positive-lag maximizer.  Hence
\begin{equation}\label{eq:strictmax}
 m\ge2,\ m\ne n \quad\Longrightarrow\quad \A_n(m)<\A_n(n).
\end{equation}
Equation \eqref{eq:strictmax} is included in the Lean development as a connector to the earlier
unique-apex theorem.

\section{Cross-prime three-atom dilution}\leavevmode\par

\begin{lemma}[Prime flanks]\label{lem:flanks}
For every integer $n\ge4$, there are primes $p,q$ with
\[
 \frac n2<p<n<q<2n.
\]
\end{lemma}

\begin{proof}
Apply Bertrand's theorem first to $\lfloor n/2\rfloor$ and then to $n$.  The weak upper endpoints
provided by Bertrand are strict here: equality would make an even integer at least four prime.
The lower inequality $n<2p$ follows from $\lfloor n/2\rfloor<p$.
\end{proof}

The next estimate uses the standard inequality $\Lam(n)\le\log n$ and antitonicity of
$x\mapsto\log x/x$ on $[e,\infty)$.

\begin{theorem}[Prime mass dominance]\label{thm:mass}
Let $n\ge4$, and let $p,q$ be as in Lemma~\ref{lem:flanks}.  Then
\begin{equation}\label{eq:masscmp}
 \w_n\le\w_p,\qquad \w_n<2\w_q.
\end{equation}
Consequently,
\begin{equation}\label{eq:twofifths}
 5\w_n<2(\w_n+\w_p+\w_q).
\end{equation}
\end{theorem}

\begin{proof}
Since $p\ge3>e$ and $p<n$, antitonicity gives
\[
 \frac{\log n}{n}\le\frac{\log p}{p}.
\]
Together with $\Lam(n)\le\log n$ and $\Lam(p)=\log p$, this proves the first inequality in
\eqref{eq:masscmp}.  Since $n<q<2n$,
\[
 \frac{\log n}{n}<\frac{2\log q}{q}.
\]
Using $\Lam(q)=\log q$ proves the second inequality.  Equation \eqref{eq:twofifths} is immediate.
\end{proof}

\begin{corollary}[Three-atom block dilution]\label{cor:primeblock}
For every $n\ge4$, primes $p,q$ may be chosen so that the apex occupies strictly less than two
fifths of the total block mass and
\begin{equation}\label{eq:blockdilute}
 \w_n\A_n(n)+\w_p\A_n(p)+\w_q\A_n(q)
 <(\w_n+\w_p+\w_q)\A_n(n).
\end{equation}
\end{corollary}

\begin{proof}
The mass assertion is \eqref{eq:twofifths}.  Theorem~\ref{thm:support} makes both prime flank
samples positive, while \eqref{eq:strictmax} makes both strictly smaller than the apex.  Their
Mangoldt masses are positive, so their weighted deficits have positive sum, proving
\eqref{eq:blockdilute}.
\end{proof}

The prime block proves strict dilution but supplies no uniform sample gap.  The next construction
removes that dependence on the position of the Bertrand primes.

\section{Sharp geometric-bracket minimax law}\leavevmode\par

The dyadic problem is an instance of a scale-free two-endpoint law.  Let $L>0$ be logarithmic
bracket width, let $\beta\ge0$ be exponential tilt, and put
\begin{equation}\label{eq:profiles}
 F_{\beta,L}(d)=e^{-\beta d}\left(1-\frac dL\right),\qquad
 G_{\beta,L}(d)=e^{\beta(L-d)}\frac dL
 \quad(0\le d\le L).
\end{equation}
Define
\begin{equation}\label{eq:balance}
 \alpha_{\beta,L}=\frac1{1+e^{\beta L}},\qquad
 C_{\beta,L}=e^{-\beta\alpha_{\beta,L}L}(1-\alpha_{\beta,L}).
\end{equation}

\begin{theorem}[Sharp tilted-tent minimax law]\label{thm:minimax}
For every $L>0$, $\beta\ge0$, and $d\in[0,L]$,
\begin{equation}\label{eq:minimax}
 \min\{F_{\beta,L}(d),G_{\beta,L}(d)\}\le C_{\beta,L}.
\end{equation}
Equality holds at $d=\alpha_{\beta,L}L$, so no smaller uniform factor can replace
$C_{\beta,L}$.  The two profiles are equal if and only if
$d=\alpha_{\beta,L}L$, and
\begin{equation}\label{eq:Cunit}
 0<C_{\beta,L}<1.
\end{equation}
\end{theorem}

\begin{proof}
Write $\alpha=\alpha_{\beta,L}$.  Direct algebra gives
\begin{equation}\label{eq:balanceidentity}
 \alpha e^{\beta L}=1-\alpha,
\end{equation}
so the profiles agree at $d=\alpha L$, with common value $C_{\beta,L}$.
Conversely, canceling the positive factor $e^{-\beta d}$ from
$F_{\beta,L}(d)=G_{\beta,L}(d)$ and using
$e^{\beta(L-d)}=e^{-\beta d}e^{\beta L}$ recovers
$1-d/L=e^{\beta L}d/L$, hence $d=\alpha L$.

The upper profile has derivative
\[
 G'_{\beta,L}(d)=\frac{e^{\beta(L-d)}}L(1-\beta d).
\]
It is positive up to the balance point.  Indeed
\[
 \beta\alpha L=\frac{\beta L}{1+e^{\beta L}}<1,
\]
where the last inequality follows from $x+1\le e^x$.  Therefore, for
$d\le\alpha L$, monotonicity gives
$\min(F(d),G(d))\le G(d)\le G(\alpha L)=C_{\beta,L}$.
For $d\ge\alpha L$, both
$e^{-\beta d}\le e^{-\beta\alpha L}$ and
$1-d/L\le1-\alpha$, with nonnegative factors, so
$\min(F(d),G(d))\le F(d)\le C_{\beta,L}$.
This proves \eqref{eq:minimax} and sharpness.  Finally $0<\alpha<1$; the exponential factor is
positive and at most one, while $0<1-\alpha<1$, proving \eqref{eq:Cunit}.
\end{proof}

The theorem is universal: it has no arithmetic input, no regularity premise on a prime sequence,
and no smallness restriction on $\beta L$.  It identifies the exact obstruction constant for any
two-endpoint geometric-bracket argument with power-law amplitude tilt.

\section{Universal subcritical tail law}\label{sec:tail}\leavevmode\par

The macroscopic construction uses a one-sided law that is stronger than pointwise contraction.
For $L>0$, $\beta\ge0$, define
\begin{equation}\label{eq:rightprofile}
 R_{\beta,L}(d)=e^{\beta d}\left(1-\frac dL\right),
 \qquad 0\le d\le L.
\end{equation}

\begin{theorem}[Sharp subcritical tail envelope]\label{thm:tail}
Assume $\beta L<1$.  Then $R_{\beta,L}$ is strictly decreasing on $[0,L]$.  For every
$0\le r\le L$ and $C\in\R$,
\begin{equation}\label{eq:tailiff}
 \bigl(\forall d\in[r,L],\ R_{\beta,L}(d)\le C\bigr)
 \quad\Longleftrightarrow\quad
 R_{\beta,L}(r)\le C.
\end{equation}
Consequently $R_{\beta,L}(r)$ is the exact least uniform ceiling on the entire tail, and
$0<d\le L$ implies $R_{\beta,L}(d)<1$.
\end{theorem}

\begin{proof}
Differentiation gives
\[
 R'_{\beta,L}(d)
 =\frac{e^{\beta d}}L\bigl(\beta(L-d)-1\bigr).
\]
On $[0,L]$, nonnegativity of $\beta$ and the subcritical inequality imply
$\beta(L-d)\le\beta L<1$.  Thus the derivative is strictly negative, and the mean-value theorem
gives strict decrease.  The forward implication in \eqref{eq:tailiff} follows by taking $d=r$;
the reverse implication follows from decrease.  The least-ceiling assertion is exactly this
equivalence.  Finally $R_{\beta,L}(0)=1$.
\end{proof}

Theorem~\ref{thm:tail} is an exact envelope statement, not merely a pointwise estimate: any
uniform ceiling on the tail must already dominate its left-endpoint value.

\section{Sharp Suzuki specialization and absolute surplus}\leavevmode\par

For $n\ge2$, put $k=\lfloor\log_2n\rfloor$, $\ell=2^k$, and $h=2^{k+1}$, so
$\ell\le n<h=2\ell$.  Set
\[
 d=\log(n/\ell)\in[0,\log2].
\]
After division by the apex value \eqref{eq:self}, the two endpoint samples are exactly
\begin{equation}\label{eq:profileidentities}
 \frac{\A_n(\ell)}{\A_n(n)}=F_{1/2,\log2}(d),\qquad
 \frac{\A_n(h)}{\A_n(n)}=G_{1/2,\log2}(d).
\end{equation}
The tilt $1/2$ is forced by the square-root amplitude in \eqref{eq:sample}; it cannot be discarded
without losing the sharp constant.

Since $e^{(\log2)/2}=\sqrt2$, the balance and contraction specialize to
\begin{equation}\label{eq:Cstar}
 \alpha_{1/2,\log2}=\sqrt2-1,
 \qquad
 C_*:=C_{1/2,\log2}
 =e^{-(\sqrt2-1)\log2/2}(2-\sqrt2).
\end{equation}

\begin{theorem}[Sharp dyadic Suzuki contraction]\label{thm:contraction}
For every integer $n\ge2$, one of $m\in\{\ell,h\}$ satisfies
\begin{equation}\label{eq:dyadicgeometry}
 n<2m,\qquad m\le2n,\qquad m\ne n,
\end{equation}
is a prime power, and obeys
\begin{equation}\label{eq:sharpcontraction}
 \A_n(m)\le C_*\A_n(n),
 \qquad
 0<C_*<\frac{\sqrt2}{2}<1.
\end{equation}
Its Mangoldt mass satisfies the absolute lower bound
\begin{equation}\label{eq:masslower}
 \w_m\ge\frac{\log2}{n}>0.
\end{equation}
\end{theorem}

\begin{proof}
Apply Theorem~\ref{thm:minimax} to \eqref{eq:profileidentities} and choose the endpoint realizing
the smaller profile.  If $n=\ell$, the lower profile equals one and the upper profile equals zero,
so the selected endpoint is $h\ne n$; otherwise both possible choices are already distinct from
$n$ when selected.  The annulus bounds follow from $\ell\le n<h=2\ell$.

Both endpoints are nontrivial powers of $2$, hence $\Lam(m)=\log2$ and
$\w_m=2\log2/m$.  Since $m\le2n$, this gives \eqref{eq:masslower}.  The formula for $C_*$ is
\eqref{eq:Cstar}.  Its exponential factor is strictly below one and
$2-\sqrt2<\sqrt2/2$, proving the strict improvement in \eqref{eq:sharpcontraction}.
\end{proof}

Define the source flank surplus
\begin{equation}\label{eq:surplusdef}
 \Sur_n(m)=\w_m\bigl(\A_n(n)-\A_n(m)\bigr).
\end{equation}

\begin{theorem}[Absolute prime-gap-free source surplus]\label{thm:surplus}
For every integer $n\ge2$, the dyadic $m$ in Theorem~\ref{thm:contraction} may be selected so that
\begin{equation}\label{eq:surplus}
 0<(1-C_*)\frac{(\log2)^2}{\sqrt n}
 \le (1-C_*)\w_m\A_n(n)
 \le \Sur_n(m).
\end{equation}
\end{theorem}

\begin{proof}
The right inequality follows by multiplying \eqref{eq:sharpcontraction} by $\w_m$ and
rearranging.  Equations \eqref{eq:masslower} and \eqref{eq:self} give
\[
 \w_m\A_n(n)\ge\frac{\log2}{n}\sqrt n\log2
 =\frac{(\log2)^2}{\sqrt n}.
\]
All factors are positive because $C_*<1$.
\end{proof}

The bound in \eqref{eq:surplus} is produced from source geometry and source mass alone.  It is not
a renamed target deficit, and it removes both prime gaps and the arithmetic value $\Lam(n)$ from
the lower bound.

\section{Exact dyadic boundary law}\leavevmode\par

Dyadic apices have a sharper identity.

\begin{theorem}[Dyadic successor boundary law]\label{thm:boundary}
For every integer $k\ge1$, put $n=2^k$ and $m=2^{k+1}$.  Then
\begin{equation}\label{eq:boundaryzero}
 \A_{2^k}(2^{k+1})=0
\end{equation}
and
\begin{equation}\label{eq:boundarysurplus}
 \Sur_{2^k}(2^{k+1})
 =\w_{2^{k+1}}\A_{2^k}(2^k)>0.
\end{equation}
\end{theorem}

\begin{proof}
The ratio $m/n$ equals $2$, so the absolute logarithmic displacement in
\eqref{eq:sample} equals $\log2$ and the tent factor vanishes.  This proves
\eqref{eq:boundaryzero}.  Substitution into \eqref{eq:surplusdef} proves the equality in
\eqref{eq:boundarysurplus}; both factors are positive.
\end{proof}

\begin{remark}[Scope of power-ray isolation]
Power-ray isolation says that distinct positive powers on one base ray cannot produce a
\emph{positive} aligned sample.  Theorem~\ref{thm:boundary} shows that this is not a no-go for block
dilution: at a dyadic apex, the adjacent same-ray atom is a positive-mass, zero-sample boundary
flank and therefore produces maximal deficit per unit of its own mass.
\end{remark}

\section{Macroscopic width-five prime-block law}\label{sec:macro}\leavevmode\par

The local dyadic surplus in \eqref{eq:surplus} tends to zero.  To reach the moving-apex scale, we
widen the support and aggregate actual primes.  Put
\begin{equation}\label{eq:fiveobjects}
 L_5=\log5,\qquad
 a_5(n)=\frac12\log(5n),\qquad
 P_n=\{p\text{ prime}:2n<p\le5n\}.
\end{equation}
Write $\A_n^{(5)}(m)$ for the positive rectangular channel of width $L_5$ and aperture $a_5(n)$,
sampled at the node $2\log m$.  Define
\begin{equation}\label{eq:qfive}
 q_5=\sqrt2\left(1-\frac{\log2}{\log5}\right).
\end{equation}

\begin{lemma}[Exact width-five aligned samples]\label{lem:fivesample}
For integers $n,m\ge5$,
\begin{equation}\label{eq:fivesample}
 \A_n^{(5)}(m)
 =\sqrt m\max\!\left\{0,\log5-\left|\log\frac mn\right|\right\}.
\end{equation}
In particular,
\begin{equation}\label{eq:fiveself}
 \A_n^{(5)}(n)=\sqrt n\log5.
\end{equation}
\end{lemma}

\begin{proof}
At the sample time $2\log m$, the fixed tent is
$2\sqrt m\max\{0,\log5-\log m\}=0$ because $m\ge5$.  The moving tent is centered at
$2a_5(n)-\log5=\log n$, and its exponential amplitude is
$e^{(2\log m)/4}=\sqrt m$.  This gives \eqref{eq:fivesample}; setting $m=n$ gives
\eqref{eq:fiveself}.
\end{proof}

\begin{theorem}[Sharp width-five flank contraction]\label{thm:fivecontraction}
One has $0<q_5<1$.  For $n\ge5$ and $p\in P_n$,
\begin{equation}\label{eq:fivecontract}
 \A_n^{(5)}(p)\le q_5\A_n^{(5)}(n).
\end{equation}
Moreover $q_5$ is the exact least uniform ceiling of the normalized continuous flank:
for every $C\in\R$,
\begin{equation}\label{eq:qsharp}
 \bigl(\forall d\in[\log2,\log5],\
 R_{1/2,\log5}(d)\le C\bigr)
 \quad\Longleftrightarrow\quad q_5\le C.
\end{equation}
\end{theorem}

\begin{proof}
For $p\in P_n$, put $d=\log(p/n)$.  Then $\log2<d\le\log5$, and
Lemma~\ref{lem:fivesample} gives
\[
 \frac{\A_n^{(5)}(p)}{\A_n^{(5)}(n)}
 =R_{1/2,\log5}(d).
\]
The certified inequality $\log5<2$ gives $(\log5)/2<1$, so
Theorem~\ref{thm:tail} bounds the ratio by its value at $\log2$.  Direct substitution gives
$R_{1/2,\log5}(\log2)=q_5$.  Strict decrease and
$0<\log2<\log5$ give $0<q_5<1$.  Equation \eqref{eq:qsharp} is
Theorem~\ref{thm:tail} with $r=\log2$.
\end{proof}

The arithmetic input is a block-mass theorem.  Let
$\vartheta(x)=\sum_{p\le x}\log p$.

\begin{lemma}[Macroscopic log-Haar mass]\label{lem:massfloor}
For every natural number $n$ with $e^{20}\le n$,
\begin{equation}\label{eq:massfloorfive}
 \sum_{p\in P_n}\frac{2\log p}{p}\ge\frac1{10}.
\end{equation}
\end{lemma}

\begin{proof}
Write $N=n$ as a real number.  The formalized Chebyshev bounds give
\begin{align}
 \vartheta(5N)-\vartheta(2N)
 &\ge (5N-1)\log2-\log(5N+2)-2\sqrt{5N}\log(5N)-2N\log4 \notag\\
 &= (N-1)\log2-\log(5N+2)-2\sqrt{5N}\log(5N).             \label{eq:thetaerror}
\end{align}
Antitonicity of $x\mapsto\log x/\sqrt x$ on $[e^2,\infty)$ and $N\ge e^{20}$ imply
\begin{equation}\label{eq:logroot}
 \frac{\log(5N)}{\sqrt{5N}}
 \le\frac{20}{e^{10}}<\frac1{50}.
\end{equation}
The strict inequality is symbolic: $e>2$, hence $e^{10}>2^{10}>1000$.  Therefore
\begin{equation}\label{eq:firsterror}
 2\sqrt{5N}\log(5N)\le\frac N5.
\end{equation}
Also $5N+2\le6N$, $\log x\le2\sqrt x$, $\sqrt6<3$, and
$\sqrt N\ge e^{10}>1000$, giving
\begin{equation}\label{eq:seconderror}
 \log(5N+2)\le6\sqrt N\le\frac N{100}.
\end{equation}
Finally $1/2<\log2<1$ and $N>20$ yield
\[
 (N-1)\log2-\frac N5-\frac N{100}
 \ge\frac{29N}{100}-\frac12\ge\frac N4.
\]
Substitution into \eqref{eq:thetaerror} proves
$\sum_{p\in P_n}\log p\ge N/4$.  Since $p\le5N$ throughout $P_n$,
\[
 \sum_{p\in P_n}\frac{2\log p}{p}
 \ge\frac2{5N}\sum_{p\in P_n}\log p
 \ge\frac1{10}.
\]
No prime table, prime-gap theorem, asymptotic remainder, or zero hypothesis enters the argument.
\end{proof}

Define the width-five source surplus
\begin{equation}\label{eq:sfivedef}
 \Sur_5(n)=\sum_{p\in P_n}\frac{2\log p}{p}
 \bigl(\A_n^{(5)}(n)-\A_n^{(5)}(p)\bigr).
\end{equation}

\begin{theorem}[Macroscopic Suzuki--Mangoldt block law]\label{thm:macro}
For every natural number $n$ with $e^{20}\le n$,
\begin{equation}\label{eq:macrolower}
 \Sur_5(n)\ge\frac{1-q_5}{10}\sqrt n\log5>0.
\end{equation}
Moreover $\Sur_5$ is unbounded: for every $B\in\R$, there is a natural number
$n\ge e^{20}$ such that
\begin{equation}\label{eq:unbounded}
 B<\Sur_5(n).
\end{equation}
\end{theorem}

\begin{proof}
Theorem~\ref{thm:fivecontraction} gives, term by term,
\[
 \A_n^{(5)}(n)-\A_n^{(5)}(p)
 \ge(1-q_5)\A_n^{(5)}(n).
\]
All prime weights are nonnegative.  Sum this inequality, apply
Lemma~\ref{lem:massfloor}, and use \eqref{eq:fiveself}; the result is
\eqref{eq:macrolower}.  Its right-hand side is positive because $q_5<1$.
For \eqref{eq:unbounded}, the coefficient
$c_5=(1-q_5)\log5/10$ is positive.  Given $B$, choose a natural number
$n>\max\{e^{20},(\max\{B,0\}/c_5)^2\}$.  Then
$B<c_5\sqrt n\le\Sur_5(n)$.
\end{proof}

Theorem~\ref{thm:macro} is the scale transition.  A single dyadic flank gives a certified surplus
of order $n^{-1/2}$, whereas the actual-prime block gives a certified surplus of order $\sqrt n$.
The gain is not a better estimate for one atom: it comes from a block carrying a fixed positive
amount of log-Haar mass while every member lies beneath one sharp flank ceiling.

\section{Formal verification}\label{sec:formal}\leavevmode\par

The universal laws are in
\begin{center}
\small\nolinkurl{GaloisForLFunctions.HTGeometry.TiltedTentGeometricBracketMinimax},\\
\small\nolinkurl{GaloisForLFunctions.HTGeometry.TiltedTentRightFlank},
\end{center}
and the Suzuki specializations are in
\begin{center}
\small\nolinkurl{GaloisForLFunctions.RHTower.Suzuki.SuzukiMangoldtApexCrossPrimeBlock},\\
\small\nolinkurl{GaloisForLFunctions.RHTower.Suzuki.SuzukiMangoldtMacroscopicPrimeBlock}.
\end{center}
The theorem correspondence is as follows.

\begin{sloppypar}
\raggedright
\begin{description}
\small
\item[Theorem~\ref{thm:support}.]
\nolinkurl{suzukiMangoldtAlignedApex_sample_pos_iff_openAnnulus}
\item[Lemma~\ref{lem:flanks}.]
\nolinkurl{exists_prime_straddling_multiplicativeAnnulus}
\item[Theorem~\ref{thm:mass}.]
\nolinkurl{mangoldtLogHaarAtomMass_le_lowerPrimeFlank};
\nolinkurl{mangoldtLogHaarAtomMass_lt_twice_upperPrimeFlank}
\item[Corollary~\ref{cor:primeblock}.]
\nolinkurl{exists_suzukiMangoldtCrossPrimeThreeAtomBlockDilution}
\item[Theorem~\ref{thm:minimax}.]
\nolinkurl{tiltedTent_geometricBracket_minimax};
\nolinkurl{tiltedTent_geometricBracket_minimax_sharp};
\nolinkurl{tiltedTent_profiles_eq_iff};
\nolinkurl{tiltedTentSharpContraction_mem_Ioo}
\item[Theorem~\ref{thm:tail}.]
\nolinkurl{tiltedTentRightProfile_strictAntiOn};
\nolinkurl{tiltedTentRightProfile_tail_bound_iff};
\nolinkurl{tiltedTentRightProfile_lt_one}
\item[Theorem~\ref{thm:contraction}.]
\nolinkurl{suzukiMangoldtDyadicSharpContraction_closedForm};
\nolinkurl{suzukiMangoldtDyadicSharpContraction_lt_sqrtTwo_div_two};
\nolinkurl{exists_dyadic_sharp_suzukiMangoldtFlankDilution}
\item[Theorem~\ref{thm:surplus}.]
\nolinkurl{suzukiMangoldt_logTwo_div_mul_alignedApex_eq};
\nolinkurl{exists_dyadic_absolute_suzukiMangoldtFlankSurplus}
\item[Theorem~\ref{thm:boundary}.]
\nolinkurl{suzukiMangoldtDyadicApex_successor_sample_eq_zero};
\nolinkurl{suzukiMangoldtDyadicApex_successor_surplus_eq_and_pos}
\item[Lemma~\ref{lem:fivesample}.]
\nolinkurl{suzukiMangoldtFiveAlignedApex_sample_eq};
\nolinkurl{suzukiMangoldtFiveAlignedApex_self_eq}
\item[Theorem~\ref{thm:fivecontraction}.]
\nolinkurl{suzukiMangoldtFiveFlankContraction_mem_Ioo};
\nolinkurl{suzukiMangoldtFiveFlankContraction_isSharp};
\nolinkurl{suzukiMangoldtFivePrimeBlock_sample_le}
\item[Lemma~\ref{lem:massfloor}.]
\nolinkurl{chebyshev_fivePrimeBlock_logMass_floor};
\nolinkurl{suzukiMangoldtFivePrimeBlock_atomMass_floor}
\item[Theorem~\ref{thm:macro}.]
\nolinkurl{suzukiMangoldtFivePrimeBlockSurplus_lower};
\nolinkurl{suzukiMangoldtFivePrimeBlockSurplus_pos};
\nolinkurl{suzukiMangoldtFivePrimeBlockSurplus_unbounded}
\end{description}

The capstone declarations depend only on Lean/mathlib's standard foundations
\texttt{propext}, \texttt{Classical.choice}, and \texttt{Quot.sound}.  The modules contain no
\texttt{sorry}, \texttt{admit}, user axiom, opaque declaration, or unsafe declaration.  No theorem
accepts RH, the external log-Haar capacity inequality, or an abstract positivity premise.
\end{sloppypar}

\section{What remains}\leavevmode\par

Theorem~\ref{thm:macro} eliminates the source-scale obstruction: the explicit block surplus is
positive, grows on the moving-apex scale, and is unbounded.  The sharpness statement
\eqref{eq:qsharp} also prevents an artificial improvement of the coefficient by reusing the same
continuous right-flank interval.

This does not complete block transport.  Remaining statement is target-side: construct disjoint
capacity-one continuum blocks for the completed positive kernel, prove their integral excess
dominates the explicit quantity in \eqref{eq:sfivedef} uniformly in $n$, and retain enough unused
kernel mass for the Gamma boundary charge.  Source positivity or source scaling can no longer be
inserted as a hypothesis; the unresolved theorem is the concrete capacity comparison itself.
Neither that comparison nor RH is assumed or claimed here.

\begin{thebibliography}{9}

\bibitem{Suzuki2026}
M.~Suzuki,
\emph{A canonical system of the Riemann zeta-function and the Riemann hypothesis},
arXiv:2606.09096, 2026.

\bibitem{RamareZuniga2026}
O.~Ramar\'e and S.~Zuniga Alterman, with an appendix by M.~Balazard,
\emph{M\"obius function and primes: an identity factory with applications},
arXiv:2312.05138v3, 2026.

\bibitem{Groskin2026}
A.~Groskin,
\emph{A finite Guinand--Weil dictionary and archimedean tail order for the truncated Weil
quadratic form},
arXiv:2607.02828, 2026.

\bibitem{MelbourneRoberto2021}
J.~Melbourne and C.~Roberto,
\emph{Transport-majorization to analytic and geometric inequalities},
arXiv:2110.03641, 2021.

\end{thebibliography}

\end{document}
